\documentclass[11pt]{article}
\usepackage[margin=1in]{geometry}
\usepackage{amsmath}
\title{Derivation Notes for Cryogenic Device Recovery Model}
\begin{document}
\maketitle

Assume the trapped-charge population decays with a first-order rate law:
\[
\frac{dQ_{\mathrm{trap}}}{dt} = -\frac{1}{\tau_{\mathrm{rec}}(T)}Q_{\mathrm{trap}}.
\]

Integrating gives
\[
Q_{\mathrm{trap}}(t,T) = Q_0 \exp\!\left(-\frac{t}{\tau_{\mathrm{rec}}(T)}\right),
\]
where $Q_0$ is the trapped-charge population immediately after irradiation or immediately after a trapping event.

Now assume the recovery time constant is thermally activated:
\[
\tau_{\mathrm{rec}}(T)=\tau_0 \exp\!\left(\frac{E_a}{k_B T}\right).
\]

Substituting into the recovery law gives
\[
Q_{\mathrm{trap}}(t,T)=Q_0 \exp\!\left[-\frac{t}{\tau_0}\exp\!\left(-\frac{E_a}{k_B T}\right)\right].
\]

This makes the low-temperature limit clear. As $T$ decreases, the factor
\[
\exp\!\left(-\frac{E_a}{k_B T}\right)
\]
becomes very small, so the effective relaxation slows strongly.

A useful engineering metric is the time to 50\% recovery, obtained from
\[
\frac{Q_{\mathrm{trap}}}{Q_0}=\frac{1}{2}.
\]
This gives
\[
t_{50}(T)=\tau_{\mathrm{rec}}(T)\ln 2.
\]

Therefore,
\[
t_{50}(T)=\tau_0 \ln 2\; \exp\!\left(\frac{E_a}{k_B T}\right),
\]
which grows rapidly as temperature decreases.

\end{document}
